% ---
% id: meat-core/thm:gradient
% kind: thm
% title: "Gradient of the pinched trace power on a face"
% summary: >-
%   The pinched trace power is differentiable at a positive definite operator,
%   and on the face of a singular one in split form: divided differences over
%   the positive spectrum, pinched operators on the ambient space. A one-sided
%   derivative is +∞ only along a direction that enlarges the support.
% prominence: standard
% status: implemented
% uses: [lem:dk, cor:concavity, foundations/fact:concavity, lem:branch-maps, lem:cone-concavity, def:branch-support, lem:branch-support]
% mentions: [meat-core/frank-wolfe/thm:firstorder, meat-core/frank-wolfe/rem:fo-nopert, meat-core/frank-wolfe/lem:fo-face, foundations/def:pinching, req:split-form-gradient, req:gradient-at-maximal-support]
% assumes: []
% refs: []
% provenance: newly-derived
% implements: certiqs_key.meat_core.objective.gradient:state_gradient
% lean:
% ---

The key-active branch operators are in general singular, and their
\tref{maximal branch support} is in general not invariant under the
\tref{key pinching} (Lemma~\ref{lem:branch-support}\,(iv)); part~(i) is therefore
stated so that no such invariance is needed, and
Requirement~\ref{req:split-form-gradient} is the rule it gives an
implementation. Part~(ii) is why gradients are taken at points of maximal
support (Requirement~\ref{req:gradient-at-maximal-support}). At maximal
support the split form gives the true one-sided derivative in every
feasible direction, because there no feasible direction enlarges the
support; that is Lemma~\ref{lem:fo-face} of the Frank--Wolfe submodule,
where Theorem~\ref{thm:firstorder} consumes it, and
Remark~\ref{rem:fo-nopert} explains why the explicit-perturbation route has
no accounted instance.

\begin{theorem}\label{thm:gradient}
Let $\alpha\in(1,\infty)$, $p=1/\alpha$, let $T$ be the pinched trace
power \eqref{eq:Tcone} on the positive semidefinite cone of a branch space
(Lemma~\ref{lem:dk}), and let $X\succ0$. Then $T$ is Fr\'echet
differentiable at $X$ with gradient
\begin{equation}\label{eq:gradT}
  \nabla T(X)
  =\alpha\,D X^{p}\Big[\cZ\big(Y^{\alpha-1}\big)\Big],
  \qquad Y\coloneqq\cZ(X^{p}),
\end{equation}
i.e.\ $\frac{\mathrm d}{\mathrm dt}T(X+tH)\big|_0=\tr[\nabla T(X)H]$,
where $DX^{p}[\cdot]$ is the self-adjoint map \eqref{eq:dk} with
$f(x)=x^{p}$ and $Y^{\alpha-1}$ is defined by functional calculus
($0^{\alpha-1}\coloneqq0$; $Y$ may be singular). Explicitly, in the
eigenbasis $X=\sum_i\lambda_iP_i$:
\begin{equation}\label{eq:gradT-kernel}
  \nabla T(X)=\alpha\sum_{i,j}
  \frac{\lambda_i^{p}-\lambda_j^{p}}{\lambda_i-\lambda_j}\,
  P_i\,\cZ\big(Y^{\alpha-1}\big)\,P_j
  \qquad(\text{diagonal kernel }p\lambda_i^{p-1}).
\end{equation}
Moreover:
\begin{enumerate}[label=(\roman*),leftmargin=2.2em]
\item \emph{(faces)} Let $X\succeq0$, write
$S\coloneqq\operatorname{supp}X$, and let $H$ be Hermitian with
$\operatorname{supp}H\subseteq S$. Then $t\mapsto T(X+tH)$ is
differentiable at $t=0$, and \eqref{eq:gradT}--\eqref{eq:gradT-kernel}
hold in the following \emph{split} form:
\begin{enumerate}[label=(\alph*),leftmargin=2.2em]
\item the divided-difference map $DX^{p}[\cdot]$ is compressed to $S$,
equivalently, the sum in \eqref{eq:gradT-kernel} runs over the strictly
positive eigenvalues of $X$ only, where $X$ is positive definite and the
kernel $f^{[1]}$ is finite;
\item $Y=\cZ(X^{p})$, $Y^{\alpha-1}$ and $\cZ(Y^{\alpha-1})$ are formed on
the \emph{ambient} branch space and are \emph{not} compressed to $S$.
\end{enumerate}
The resulting $\nabla T(X)$ determines the derivative in every direction
supported in $S$; its components outside $S$ are neither determined by
this nor needed, support-enlarging directions being the subject of (ii).
Compressing the \emph{whole} computation to a subspace $S'$ is licensed
when $\sum_{j\in\Zset}Z_jS\subseteq S'$, in particular when $S\subseteq S'$
and $S'$ is $\cZ$-invariant ($Z_jS'\subseteq S'$ for every $j\in\Zset$):
then the ambient complement contributes $0$ to $T$, and the compressed and
ambient computations agree. Compressing to $S$ itself is licensed
\emph{only} when $S$ is $\cZ$-invariant; otherwise $\cZ$ transports weight
out of $S$ and the compressed value of $T$ is strictly smaller (the proof
gives the argument and a rank-one witness). For the branch operators $X=\cG_{\hat c,s}(\omega_s)$, whose
faces are in general not $\cZ$-invariant, the relevant face is
$S=W_{\hat c,s}$ of \eqref{eq:Wcs} at maximal support, and (a)--(b) is
what makes the gradient available there.
\item \emph{(one-sided derivatives)} For $X\succeq0$ and any Hermitian
$H$ with $X+tH\succeq0$ for small $t>0$, the one-sided directional
derivative $\lim_{t\downarrow0}\frac{T(X+tH)-T(X)}{t}$ exists in
$(-\infty,+\infty]$, and can equal $+\infty$ only when $H$ strictly
enlarges the support of $X$.
\end{enumerate}
\end{theorem}

\begin{proof}
\emph{Idea.} At $X\succ0$ the gradient is the chain rule through
$X\mapsto X^{p}$ (Lemma~\ref{lem:dk}), the linear pinching and
$Y\mapsto\tr[Y^{\alpha}]$. On a face the inner map is differentiated
inside the support, where $X$ is positive definite, while the outer map is
differentiable at every $Y\succeq0$ and so stays on the ambient space.
Concavity of $T$ locates the only infinite one-sided derivatives.

\emph{Step 1 (differentiability at $X\succ0$).} $X\mapsto X^{p}$ is
Fr\'echet differentiable at $X\succ0$ (Lemma~\ref{lem:dk} with
$f=x^{p}$ on $I=(0,\infty)$); $\cZ$ is linear; and with
$g(y)=\max\{y,0\}^{\alpha}$ of Lemma~\ref{lem:dk}, continuously
differentiable on $\mathbb{R}$ because $\alpha>1$, the map
$Y\mapsto\tr g(Y)$, which is $\tr[Y^{\alpha}]$ on the positive
semidefinite operators, is Fr\'echet differentiable at every $Y\succeq0$
with derivative $K\mapsto\alpha\tr[Y^{\alpha-1}K]$ by
\eqref{eq:trace-diff}. For $|t|$ small $X+tH\succ0$, so
$T(X+tH)=\tr g\big(\cZ((X+tH)^{p})\big)$. The chain rule gives
\begin{equation*}
  \frac{\mathrm d}{\mathrm dt}T(X+tH)\Big|_0
  =\alpha\tr\big[Y^{\alpha-1}\,\cZ\big(DX^{p}[H]\big)\big]
  =\alpha\tr\big[\cZ\big(Y^{\alpha-1}\big)\,DX^{p}[H]\big]
  =\tr\big[\alpha\,DX^{p}\big[\cZ(Y^{\alpha-1})\big]\,H\big],
\end{equation*}
using self-adjointness of the pinching and then of $DX^{p}$
(Lemma~\ref{lem:dk}). \eqref{eq:gradT-kernel} is \eqref{eq:dk} written
out.

\emph{Step 2 (faces, part (i)).} Write $S=\operatorname{supp}X$ and let
$\Pi_{S}$ be its projector. For Hermitian $H$ with
$\operatorname{supp}H\subseteq S$ we have
$X+tH=\Pi_{S}(X+tH)\Pi_{S}$, and $X$ is positive definite on $S$, so for
$|t|$ small the compression $(X+tH)|_{S}$ is positive definite and
$(X+tH)^{p}=\big((X+tH)|_{S}\big)^{p}\oplus0$ by the convention $0^{p}=0$
of \xref{foundations/def:pinching}. Lemma~\ref{lem:dk}, applied inside
$\Lin{S}$ with $f(x)=x^{p}$ on $I=(0,\infty)$, which contains
$\operatorname{spec}(X|_{S})$, gives
\begin{equation}\label{eq:face-dpow}
  \frac{\mathrm d}{\mathrm dt}(X+tH)^{p}\Big|_{0}
  =\!\!\sum_{i,j:\ \lambda_i,\lambda_j>0}\!\!
  \frac{\lambda_i^{p}-\lambda_j^{p}}{\lambda_i-\lambda_j}\,P_iHP_j
  \eqqcolon D_{S}X^{p}[H],
\end{equation}
a self-adjoint map on the Hermitian operators supported in $S$. No term
with $\lambda_i=0$ or $\lambda_j=0$ occurs, which removes the divergence
of the kernel at the origin; this is (a). The outer factor is untouched by
that compression: with $g(y)=\max\{y,0\}^{\alpha}$ as in Step~1,
\eqref{eq:trace-diff} makes $Y\mapsto\tr g(Y)$, which is $\tr[Y^{\alpha}]$
on the positive semidefinite operators, differentiable at every
$Y\succeq0$, in particular at the \emph{ambient}
$Y=\cZ(X^{p})$, which is not supported in $S$ unless $S$ is
$\cZ$-invariant. This is (b). Composing the two,
\begin{equation*}
\begin{aligned}
  \frac{\mathrm d}{\mathrm dt}T(X+tH)\Big|_{0}
  &=\alpha\tr\big[Y^{\alpha-1}\,\cZ\big(D_{S}X^{p}[H]\big)\big]
  =\alpha\tr\big[\cZ\big(Y^{\alpha-1}\big)\,D_{S}X^{p}[H]\big]\\
  &=\tr\big[\alpha\,D_{S}X^{p}\big[\cZ(Y^{\alpha-1})\big]\,H\big],
\end{aligned}
\end{equation*}
by self-adjointness of the pinching and then of $D_{S}X^{p}$
(Lemma~\ref{lem:dk} inside $\Lin{S}$; only
$\Pi_{S}\cZ(Y^{\alpha-1})\Pi_{S}$ enters, because $H$ is supported in
$S$).

\emph{Step 3 (compression).} Put $S_{\cZ}\coloneqq\sum_{j\in\Zset}Z_jS$.
It contains $S$, because $\sum_jZ_j=\idon{}$, and it is $\cZ$-invariant,
because $Z_iZ_j=\delta_{ij}Z_i$; and for every $W$ supported in $S_{\cZ}$,
$\cZ(W)=\sum_jZ_jWZ_j$ is supported in $\sum_jZ_jS_{\cZ}=S_{\cZ}$. With
$X^{p}$ supported in $S$, therefore, $Y=\cZ(X^{p})$, $Y^{\alpha-1}$ and
$\cZ(Y^{\alpha-1})$ are supported in $S_{\cZ}$. Let
$S_{\cZ}\subseteq S'$, which holds in particular when $S\subseteq S'$ and
$Z_jS'\subseteq S'$ for all $j$. Compressing the pinching to $S'$,
$W\mapsto\Pi_{S'}\cZ(W)\Pi_{S'}$, then changes none of these operators, so
every trace above may be taken inside $S'$ and the ambient complement
contributes $0$: the compressed and ambient computations agree.

If $S$ is not $\cZ$-invariant, compressing to $S$ itself changes $T$
already. Then $S_{\cZ}\ne S$, and $\operatorname{supp}Y=S_{\cZ}$: a vector
$\ket{w}$ with $\bra{w}Z_jX^{p}Z_j\ket{w}=0$ for every $j$ has
$X^{p/2}Z_j\ket{w}=0$, hence $Z_j\ket{w}\perp S$, for every $j$, and
conversely each $Z_jX^{p}Z_j$ has range in $Z_jS$. So $Y$ is not supported
in $S$, and the compressed value $\tr[(\Pi_SY\Pi_S)^{\alpha}]$ is strictly
smaller than $T(X)=\tr[Y^{\alpha}]$: for a unit eigenvector $u\in S$ of
$\Pi_SY\Pi_S$, with eigenvalue $\mu=\bra{u}Y\ket{u}$, Jensen's inequality
for the convex $x\mapsto x^{\alpha}$ and the spectral distribution of $Y$
in $u$ gives $\mu^{\alpha}\le\bra{u}Y^{\alpha}\ket{u}$; summing over an
orthonormal eigenbasis of $S$ gives
$\tr[(\Pi_SY\Pi_S)^{\alpha}]\le\tr[\Pi_SY^{\alpha}]$, the operator
$\Pi_SY\Pi_S$ vanishing on $S^{\perp}$; and
$\tr[\Pi_SY^{\alpha}]<\tr[Y^{\alpha}]$, because $Y^{\alpha}\succeq0$ has the
support of $Y$, which is not contained in $S$. Witness: let
$\ket\eta=a\ket{j}_R\ket{u}+b\ket{j'}_R\ket{u'}$ with $j\ne j'$,
$|a|^{2}+|b|^{2}=1$, and $X=\dyad{\eta}$, so that $X^{p}=X$ and
$\cZ(X^{p})=|a|^{2}\dyad{j,u}+|b|^{2}\dyad{j',u'}$ has rank two while $X$
has rank one. The ambient value is
$T(X)=|a|^{2\alpha}+|b|^{2\alpha}$, whereas compressing to
$S=\mathbb{C}\ket\eta$ would return
$\big(\bra\eta\cZ(X^{p})\ket\eta\big)^{\alpha}
=(|a|^{4}+|b|^{4})^{\alpha}$; at $|a|^{2}=|b|^{2}=\tfrac12$ these are
$2^{1-\alpha}$ and $2^{-\alpha}$.

\emph{Step 4 (branch operators).} For $X=\cG_{\hat c,s}(\omega_s)$,
\eqref{eq:Wcs-bound} gives $S\subseteq W_{\hat c,s}$, with equality at
maximal support, and $W_{\hat c,s}$ need not be $\cZ$-invariant
(Lemma~\ref{lem:branch-support}\,(iv)). The split form (a)--(b) is
therefore the operative one: it needs no $\cZ$-invariance, only
$\operatorname{supp}H\subseteq S$. The whole computation may be
compressed to a $\cZ$-invariant $S'\supseteq W_{\hat c,s}$, for instance
$\sum_{j\in\Zset}Z_jW_{\hat c,s}$, which is $\cZ$-invariant because
$Z_iZ_j=\delta_{ij}Z_i$.

\emph{Step 5 (one-sided derivatives, part (ii)).} $T$ is concave on the
positive semidefinite cone by Lemma~\ref{lem:cone-concavity}\,(b), which
carries the concavity of \xref{foundations/fact:concavity}, in the form of
Corollary~\ref{cor:concavity} on $\Sleq{}$, to the whole cone by the
positive $1$-homogeneity $T(cX)=c\,T(X)$ for $c\ge0$. For a concave
function, $t\mapsto\frac{T(X+tH)-T(X)}{t}$ is non-increasing on the
feasible ray, so the limit exists in $(-\infty,+\infty]$. If $H$
preserves the support, (i) gives a finite value. The scalar example
$x\mapsto x^{p}$ at $x=0$ shows the support-enlarging derivative can be
$+\infty$.
\end{proof}
